% =====================================================================
%  bg_C_new_sections.tex
%  WHERE:  at the END of Chapter 2 "Background", i.e. after everything
%          from bg_B_language_models.tex and BEFORE \chapter{Related Work}.
%  ADDS:   2.4 Neuro-Symbolic Artificial Intelligence
%          2.5 Knowledge-Enhanced Large Language Models
%          2.6 Automatic Prompt Optimization
%          2.7 Cardiovascular Disease and Patient Education
%          2.8 Summary
% =====================================================================

% ---------------------------------------------------------------------
\section{Neuro-Symbolic Artificial Intelligence}
\label{sec:bg-nesy}
% ---------------------------------------------------------------------

Artificial intelligence has developed along two traditions. \emph{Symbolic}
approaches represent knowledge explicitly, for example as ontologies, logical
rules or Knowledge Graphs, and derive conclusions by transparent reasoning.
They are interpretable and precise, but they are expensive to build and
struggle with noisy, unstructured input such as free text.
\emph{Neural} (sub-symbolic) approaches, including LLMs, learn statistical
patterns from large amounts of data. They are robust to noise and very good
at producing language, but their knowledge is stored implicitly in millions
of parameters and their conclusions are difficult to verify.

\emph{Neuro-symbolic AI} (NeSy) aims to combine both traditions so that
each compensates for the weaknesses of the other~\cite{garcez2023neurosymbolic,kautz2022third}.
Kautz~\cite{kautz2022third} distinguishes several integration patterns,
ranging from neural models that call a symbolic solver, to symbolic knowledge
that is \emph{compiled into} the training data or the loss function of a
neural network. This thesis follows the latter pattern: symbolic rules learnt
from a Knowledge Graph are translated into natural-language training
examples, so that the neural model internalises the explicit knowledge
during fine-tuning.

\subsection{Rule Mining over Knowledge Graphs}
\label{sec:bg-rule-mining}

A Knowledge Graph contains many regularities that are not stated as rules
but can be learnt from the data, for example ``if two patients answered the
same question, then \ldots''. \emph{Rule mining} discovers such regularities
automatically and expresses them as logical rules.

\begin{bgdefinition}{Horn Rule}
An \emph{atom} is an expression $p(X,Y)$ where $p \in L$ is an edge label
(predicate) and $X, Y$ are variables or constants. A \emph{Horn rule} has the
form
\[
  \underbrace{B_1 \wedge B_2 \wedge \dots \wedge B_n}_{\text{body}}
  \;\Rightarrow\;
  \underbrace{r(x,y)}_{\text{head}},
\]
where the body is a conjunction of atoms and the head is a single atom. A rule
is \emph{closed} if every variable appears at least twice, which ensures
that it makes concrete predictions. Replacing all variables by constants such
that every body atom is a triple in $E$ is called an \emph{instantiation} of
the rule; the instantiated head is a \emph{prediction}~\cite{galarraga2013amie,galarraga2015fast}.
\end{bgdefinition}

Not every rule that can be formed is useful, so candidate rules are ranked
by quality measures. Under the Open-World Assumption, a prediction that is
missing from the graph cannot simply be counted as an error
(Section~\ref{sec:bg-owa}). The AMIE family of rule miners therefore
introduced the \emph{Partial Completeness Assumption} (PCA): if the graph
knows \emph{at least one} object $y'$ for a subject $x$ and relation $r$,
it is assumed to know \emph{all} of them; only in this case is a missing
prediction $r(x,y)$ treated as a counter-example~\cite{galarraga2013amie}.

\begin{bgdefinition}{Rule Quality Measures}
Let $\vec{B} \Rightarrow r(x,y)$ be a closed Horn rule and let
$\vec{z}$ denote the variables of the body other than $x$ and $y$.
\begin{align*}
  \mathit{supp}(\vec{B} \Rightarrow r(x,y)) &=
     \#\,(x,y) : \exists \vec{z} : \vec{B} \wedge r(x,y)\\[2pt]
  \mathit{hc}(\vec{B} \Rightarrow r(x,y)) &=
     \frac{\mathit{supp}(\vec{B} \Rightarrow r(x,y))}{\#\,(x',y') : r(x',y')}\\[2pt]
  \mathit{conf}(\vec{B} \Rightarrow r(x,y)) &=
     \frac{\mathit{supp}(\vec{B} \Rightarrow r(x,y))}{\#\,(x,y) : \exists \vec{z} : \vec{B}}\\[2pt]
  \mathit{conf}_{\mathit{PCA}}(\vec{B} \Rightarrow r(x,y)) &=
     \frac{\mathit{supp}(\vec{B} \Rightarrow r(x,y))}{\#\,(x,y) : \exists \vec{z}, y' : \vec{B} \wedge r(x,y')}
\end{align*}
The \emph{support} counts correct predictions, the \emph{head coverage}
(\textit{hc}) normalises the support by the size of the head relation, the
\emph{standard confidence} treats every missing prediction as false, and the
\emph{PCA confidence} only counts predictions as false if the graph knows
another object for the same subject~\cite{galarraga2013amie,galarraga2015fast}.
\end{bgdefinition}

AMIE~3~\cite{lajus2020fast} is the current version of this rule miner. It
explores the space of closed Horn rules with a refinement strategy, prunes
rules whose head coverage falls below a threshold, and computes all measures
exactly while scaling to graphs with millions of facts. Figure~\ref{fig:bg-rule-mining}
illustrates how a rule is supported by a Knowledge Graph and how it predicts
a missing fact.

\begin{figure}[htbp]
  \centering
  \begin{tikzpicture}[
      font=\small, >={Stealth[length=2.2mm]},
      dis/.style={draw=bgBlue, thick, fill=bgBlueLight, rounded corners=3pt,
                  minimum width=3.1cm, minimum height=0.8cm, align=center},
      drug/.style={draw=bgGreen, thick, fill=bgGreenLight, ellipse,
                   minimum width=1.9cm, minimum height=0.8cm},
      edge/.style={->, thick, bgGrey},
      pred/.style={->, very thick, dashed, bgOrange},
      lab/.style={font=\scriptsize, fill=white, inner sep=1pt}]
    \node[dis] (chf) at (0,3.2) {Congestive\\heart failure};
    \node[dis] (ah)  at (0,-0.4) {Arterial\\hypertension};
    \node[drug] (d1) at (-4.2,1.4) {Drug 1};
    \node[drug] (d2) at (-1.4,1.4) {Drug 2};
    \node[drug] (d3) at (1.4,1.4)  {Drug 3};
    \node[drug] (d4) at (4.2,1.4)  {Drug 4};
    \foreach \d in {d1,d2,d3,d4} \draw[edge] (chf) -- (\d);
    \draw[edge] (ah) -- (d1);
    \draw[edge] (ah) -- (d2);
    \draw[edge] (ah) -- (d3);
    \draw[pred] (ah) -- (d4) node[midway, below right, lab, text=bgOrange] {predicted};
    % legend / rule
    \node[draw=bgGrey, rounded corners=3pt, fill=bgGreyLight, align=left,
          font=\footnotesize, text width=12.2cm, anchor=north] at (0,-1.35)
      {\textbf{Rule:}\;
       $\mathit{treatedBy}(\textit{CHF}, b) \;\Rightarrow\; \mathit{treatedBy}(\textit{AH}, b)$\\[2pt]
       Solid edges: facts in the graph (\textit{treatedBy}). Drugs~1--3 satisfy
       body \emph{and} head (support $=3$); Drug~4 satisfies only the body.\\
       $\mathit{conf} = 3/4 = 0.75$. The dashed edge is the rule's
       prediction for the missing fact.};
  \end{tikzpicture}
  \caption[Rule mining over a Knowledge Graph]{Rule mining over a toy
  Knowledge Graph. The rule is learnt from drugs that treat both diseases and
  is then used to predict a missing link (dashed). \textit{treatedBy}
  abbreviates the Wikidata-derived relation
  \textit{drug\_or\_therapy\_used\_for\_treatment}.}
  \label{fig:bg-rule-mining}
\end{figure}

\begin{bgexample}[A rule mined from the guideline Knowledge Graph]
Running AMIE~3 on the guideline Knowledge Graph of this thesis (enriched with
Wikidata) produces, among others, the rule
\begin{align*}
  &\mathit{drug\_or\_therapy\_used\_for\_treatment}(\textit{congestive\_heart\_failure},\, b)\\
  \Rightarrow\; &\mathit{drug\_or\_therapy\_used\_for\_treatment}(\textit{arterial\_hypertension},\, b)
\end{align*}
with $\mathit{supp}=52$, a body size of $65$, and therefore
$\mathit{conf} = \mathit{conf}_{\mathit{PCA}} = 52/65 = 0.80$; its head
coverage is $0.032$. In words: \emph{80\,\% of the drugs that the graph
lists for congestive heart failure are also listed for arterial
hypertension.} This reflects a real pharmacological overlap (e.g.\
diuretics, ACE inhibitors and beta-blockers are used for both conditions).
However, a mined rule is a \emph{statistical regularity of the graph}, not a
clinical recommendation: the remaining 20\,\% of drugs show that the rule
does not hold in general, and its predictions must not be presented to
patients without verification.
\end{bgexample}

\subsection{From Symbolic Rules to Chain-of-Thought Training Data}
\label{sec:bg-rules-to-cot}

Mined rules are formulated in the vocabulary of the Knowledge Graph and
cannot be used directly to train a language model. They are therefore
\emph{grounded} in concrete entities and \emph{verbalised} as natural
language.

\begin{bgdefinition}{Rule Grounding}
Let $\rho: \vec{B} \Rightarrow r(x,y)$ be a Horn rule. A \emph{grounding} of
$\rho$ is a substitution $\theta$ that maps every variable of $\rho$ to a
constant in $N$ such that every grounded body atom $\theta(B_i)$ is a
triple in $E$. The grounded head $\theta(r(x,y))$ is the conclusion
supported by the grounded body, which serves as its evidence or
\emph{reasoning path}~\cite{galarraga2015fast,zhou2026breaks}.
\end{bgdefinition}

A grounding provides exactly the ingredients of a chain-of-thought
demonstration (Definition of Chain-of-Thought in
Section~\ref{sec:bg-prompting}): the grounded body atoms are the
intermediate reasoning steps, the rule is the justification that connects
them, and the grounded head is the answer. Rules with low PCA confidence can
be excluded or presented as uncertain, so that the model only learns
reasoning patterns that are well supported by the data.

\begin{bgexample}[From a rule to a chain-of-thought training instance (simplified)]
\textbf{Rule} (guideline KG, $\mathit{conf}_{\mathit{PCA}} = 1.0$):
\begin{align*}
  &\mathit{gradingSystemName}(a, \textit{``ESC class of recommendation''})\\
  \Rightarrow\; &\mathit{gradeDimension}(a, \textit{CLASS\_OF\_RECOMMENDATION})
\end{align*}
\textbf{Grounding:} $a \mapsto g_{17}$, a grade record attached to one
guideline recommendation.\par\smallskip
\textbf{Generated training instance:}\\
\emph{Question:} Which grade dimension does grade record $g_{17}$ describe?\\
\emph{Reasoning:} (1)~The knowledge graph states that $g_{17}$ uses the
grading system ``ESC class of recommendation''. (2)~Records graded with this
system describe the dimension \emph{class of recommendation} (rule
confidence 1.0). (3)~Therefore, $g_{17}$ describes a class of
recommendation.\\
\emph{Answer:} CLASS\_OF\_RECOMMENDATION
\end{bgexample}

% ---------------------------------------------------------------------
\section{Knowledge-Enhanced Large Language Models}
\label{sec:bg-kg-llm}
% ---------------------------------------------------------------------

Knowledge Graphs and LLMs have complementary strengths: Knowledge Graphs are
explicit, verifiable and easy to update, while LLMs are flexible and
produce fluent language. Pan et al.~\cite{pan2024unifying} distinguish three
ways of combining them: (i)~\emph{KG-enhanced LLMs}, which use a Knowledge
Graph during pre-training, fine-tuning or inference to improve the factual
quality of the model; (ii)~\emph{LLM-augmented KGs}, which use LLMs to build,
complete or query Knowledge Graphs; and (iii)~\emph{synergised} approaches,
in which both components improve each other. This thesis belongs mainly to
the first category: rules from the Knowledge Graph shape the fine-tuning data,
and facts from the Knowledge Graph are given to the model as context at
inference time.

\subsection{Retrieval-Augmented Generation}
\label{sec:bg-rag}

The most common way of providing knowledge at inference time is
\emph{Retrieval-Augmented Generation} (RAG)~\cite{lewis2020retrieval}.

\begin{bgdefinition}{Retrieval-Augmented Generation}
Given an input $x$ and a knowledge source $\mathcal{D}$, a retriever selects
a set of relevant items $Z = \{z_1, \dots, z_k\} \subseteq \mathcal{D}$, and
a generator produces the output conditioned on both:
\[
  p(y \mid x) \;\approx\; \sum_{z \in Z} p_{\eta}(z \mid x)\; p_{\theta}(y \mid x, z),
\]
where $p_{\eta}$ is the retrieval model and $p_{\theta}$ the language
model~\cite{lewis2020retrieval}. In \emph{KG-RAG}, $\mathcal{D}$ is a
Knowledge Graph and the retrieved items are triples, paths or sub-graphs that
are serialised as text before being added to the prompt.
\end{bgdefinition}

Because the knowledge is kept outside of the model, RAG allows updates
without re-training, for example when a new version of a guideline is
published, and the retrieved items can be shown to the user as evidence.
However, KG-RAG implicitly assumes that the answer can be found in the graph.
As discussed in Section~\ref{sec:bg-owa}, this is rarely true. Zhou et
al.~\cite{zhou2026breaks} therefore proposed BRINK, a benchmark in which the
triple that directly answers a question is \emph{removed}, while an
alternative reasoning path remains. A system can only answer correctly if it
\emph{reasons} over the remaining facts instead of simply looking up the
answer (Figure~\ref{fig:bg-kgrag}).

\begin{figure}[htbp]
  \centering
  \begin{tikzpicture}[
      font=\small, >={Stealth[length=2.2mm]},
      box/.style={draw, thick, rounded corners=3pt, minimum height=1.1cm,
                  align=center, text width=2.4cm},
      ent/.style={draw=bgBlue, thick, fill=white, circle, minimum size=0.75cm,
                  inner sep=0pt, font=\footnotesize},
      lab/.style={font=\scriptsize, inner sep=1.5pt}]
    % top row
    \node[box, fill=bgGreyLight] (q) at (0,0) {\textbf{Question}\\\footnotesize ``Which grade\\dimension does $g$\\describe?''};
    \node[box, fill=bgPurpleLight, draw=bgPurple] (r) at (3.75,0) {\textbf{Retriever}};
    \node[box, fill=bgOrangeLight, draw=bgOrange] (llm) at (8.0,0) {\textbf{LLM}\\\footnotesize reasons over\\the evidence};
    \node[box, fill=bgGreenLight, draw=bgGreen] (ans) at (11.6,0) {\textbf{Answer} $c$\\\footnotesize + evidence\\path};
    % incomplete KG below the retriever
    \begin{scope}[shift={(3.75,-3.6)}]
      \node[ent] (g) at (-1.9,0.55) {$g$};
      \node[ent] (s) at (1.9,0.55)  {$s$};
      \node[ent] (c) at (0,-1.05)   {$c$};
      \draw[->, thick] (g) -- node[lab, above] {gradingSystem} (s);
      \draw[->, thick] (s) -- node[lab, below right] {describes} (c);
      \draw[->, thick, dashed, red!70!black] (g) -- node[lab, below left, text=red!70!black] {removed} (c);
    \end{scope}
    \begin{scope}[on background layer]
      \node[draw=bgBlue, fill=bgBlueLight, rounded corners=6pt, fit=(g)(s)(c),
            inner xsep=22pt, inner ysep=12pt] (kg) {};
    \end{scope}
    \node[font=\footnotesize\bfseries, text=bgBlue, anchor=north east] at (kg.south east) {Incomplete KG\ };
    % arrows
    \draw[->, thick] (q) -- (r);
    \draw[->, thick] ([xshift=-4mm]r.south) -- node[lab, left] {query} ([xshift=-4mm]r.south |- kg.north);
    \draw[<-, thick] ([xshift=4mm]r.south) -- node[lab, right, align=left] {path\\$g \rightarrow s \rightarrow c$} ([xshift=4mm]r.south |- kg.north);
    \draw[->, thick] (r) -- node[lab, above, align=center] {question\\+ path} (llm);
    \draw[->, thick] (llm) -- (ans);
  \end{tikzpicture}
  \caption[KG-based Retrieval-Augmented Generation under incomplete knowledge]{
  KG-based Retrieval-Augmented Generation under incomplete knowledge. The
  directly supporting triple ($g \rightarrow c$, red) is missing, so the LLM
  must infer the answer from the alternative path $g \rightarrow s \rightarrow c$.
  This is the evaluation setting of the BRINK benchmark~\cite{zhou2026breaks}.}
  \label{fig:bg-kgrag}
\end{figure}

\subsection{Evaluating Answers against Answer Sets}
\label{sec:bg-metrics}

Many questions over a Knowledge Graph have more than one correct answer
(e.g.\ ``Which drugs are recommended for HFrEF?''). Model outputs are
therefore compared with a \emph{set} of gold answers rather than with a single
string.

\begin{bgdefinition}{Answer-Set Metrics}
For a question $q$, let $A_q$ be the set of gold answers and $P_q$ the set
of answers extracted from the model output after normalisation. Then
\[
  \mathit{Hits@Any}(q) = \mathbb{1}\left[P_q \cap A_q \neq \emptyset\right], \quad
  \mathit{Prec}(q) = \frac{|P_q \cap A_q|}{|P_q|}, \quad
  \mathit{Rec}(q) = \frac{|P_q \cap A_q|}{|A_q|},
\]
\[
  F_1(q) = \frac{2\cdot \mathit{Prec}(q)\cdot \mathit{Rec}(q)}{\mathit{Prec}(q)+\mathit{Rec}(q)} .
\]
\emph{Hits@Hard} additionally checks whether $P_q$ contains the specific
\emph{hard} answer whose supporting triple was removed. All metrics are
averaged over the questions of a test set~\cite{zhou2026breaks}.
\end{bgdefinition}

\begin{bgexample}[Computing the answer-set metrics]
Let $A_q = \{\text{ACE inhibitor}, \text{ARNI}, \text{beta-blocker}, \text{MRA}\}$
and suppose the model answers ``beta-blocker, MRA, digoxin'', i.e.\
$P_q = \{\text{beta-blocker}, \text{MRA}, \text{digoxin}\}$. Then
$|P_q \cap A_q| = 2$, so $\mathit{Hits@Any}=1$,
$\mathit{Prec} = 2/3 \approx 0.67$, $\mathit{Rec} = 2/4 = 0.50$ and
$F_1 = (2 \cdot 0.67 \cdot 0.50)/(0.67+0.50) \approx 0.57$. The model found
some correct answers (high Hits@Any), but it missed half of them and added
one that is not in the gold set.
\end{bgexample}

% ---------------------------------------------------------------------
\section{Automatic Prompt Optimization}
\label{sec:bg-apo}
% ---------------------------------------------------------------------

The output of an LLM depends strongly on the wording of its prompt
(Section~\ref{sec:bg-prompting}). Writing good prompts by hand is a
trial-and-error process that does not scale and is hard to reproduce.
\emph{Automatic prompt optimization} treats the prompt itself as the object
that is optimised, while the model stays fixed.

\begin{bgdefinition}{Prompt Optimization Problem}
Let $\mathcal{M}$ be a fixed language model, $\mathcal{P}$ the space of
possible prompts, $D_{\mathit{val}}$ a set of examples $(x, y^{*})$ and
$f(\cdot) \in [0,1]$ a task-specific scoring function. Prompt optimization
searches for
\[
  p^{*} \;=\; \arg\max_{p \in \mathcal{P}} \;
  \frac{1}{|D_{\mathit{val}}|} \sum_{(x,y^{*}) \in D_{\mathit{val}}}
  f\big(\mathcal{M}(p, x),\, y^{*}\big).
\]
Because $\mathcal{P}$ is discrete and $\mathcal{M}$ is used as a black box,
gradient-based optimisation is not possible, and search heuristics are
used instead~\cite{pryzant2023automatic,yang2024large,secheresse2025gaapo}.
\end{bgdefinition}

Several strategies for this search have been proposed:

\begin{itemize}
  \item \textbf{Automatic Prompt Optimization (APO/ProTeGi)}~\cite{pryzant2023automatic}
        evaluates a prompt on a mini-batch, asks an LLM to write a
        natural-language critique of the errors (a ``textual gradient''), and
        then asks it to edit the prompt in the opposite direction of this
        critique.
  \item \textbf{Optimization by PROmpting (OPRO)}~\cite{yang2024large} shows an
        optimiser LLM a list of previous prompts together with their scores
        and asks it to propose a new prompt that is expected to score higher.
  \item \textbf{Genetic algorithms}~\cite{holland1975adaptation,goldberg1989genetic}
        maintain a \emph{population} of candidate solutions and improve it over
        \emph{generations} by \emph{mutation} (random changes),
        \emph{crossover} (combining two parents) and \emph{selection}
        (keeping the fittest candidates).
\end{itemize}

\emph{GAAPO} (Genetic Algorithmic Applied to Prompt
Optimization)~\cite{secheresse2025gaapo} combines these ideas. It is a
genetic algorithm whose new candidates are produced by several different
generators: LLM-based mutation with different strategies (e.g.\ adding an
expert persona, decomposing the task, making the prompt more concise),
crossover of two parent prompts, few-shot augmentation with training
examples, APO and OPRO. All candidates are scored on the same validation
set, and the best ones are kept as parents for the next generation
(Figure~\ref{fig:bg-gaapo}). Since every candidate records which generator
produced it, GAAPO also reveals which optimisation strategy works best for a
given task.

\begin{figure}[htbp]
  \centering
  \begin{tikzpicture}[
      font=\small, >={Stealth[length=2.2mm]},
      stage/.style={draw, thick, rounded corners=4pt, minimum height=1.1cm,
                    text width=2.9cm, align=center},
      op/.style={draw=bgOrange, fill=bgOrangeLight, rounded corners=2pt,
                 minimum width=2.3cm, minimum height=0.5cm, font=\footnotesize}]
    \node[stage, fill=bgGreyLight] (init) at (0,0) {\textbf{Initial prompt(s)}\\\footnotesize population $P_0$};
    \node[stage, fill=bgBlueLight, draw=bgBlue] (pop) at (0,-2.4) {\textbf{Population} $P_{g-1}$\\\footnotesize best $N$ prompts};
    % operators
    \node[op] (o1) at (5.0,0.2)  {Mutation};
    \node[op] (o2) at (5.0,-0.5) {Crossover};
    \node[op] (o3) at (5.0,-1.2) {Few-shot};
    \node[op] (o4) at (5.0,-1.9) {APO};
    \node[op] (o5) at (5.0,-2.6) {OPRO};
    \node[draw=bgOrange, dashed, rounded corners=4pt, fit=(o1)(o5), inner sep=6pt,
          label={[font=\footnotesize\bfseries, text=bgOrange]above:1. Generation}] (gen) {};
    \node[stage, fill=bgPurpleLight, draw=bgPurple] (eval) at (5.0,-5.0)
        {\textbf{2. Evaluation}\\\footnotesize score $f$ on $D_{\mathit{val}}$\\with fixed LLM $\mathcal{M}$};
    \node[stage, fill=bgGreenLight, draw=bgGreen] (sel) at (0,-5.0)
        {\textbf{3. Selection}\\\footnotesize keep top-$N$\\(elitism)};
    \node[stage, fill=white] (best) at (-4.2,-5.0) {\textbf{Best prompt} $p^{*}$\\\footnotesize after $G$ generations};
    \draw[->, thick] (init) -- (pop);
    \draw[->, thick] (pop.east) -- node[above, font=\scriptsize] {parents} (gen.west|-pop.east);
    \draw[->, thick] (gen.south) -- node[right, font=\scriptsize] {new candidates} (eval.north);
    \draw[->, thick] (eval) -- (sel);
    \draw[->, thick] (sel) -- node[left, font=\scriptsize, align=right] {next\\generation} (pop);
    \draw[->, thick] (sel) -- node[above, font=\scriptsize] {$g = G$} (best);
  \end{tikzpicture}
  \caption[The GAAPO optimisation loop]{The GAAPO optimisation loop. In each
  generation, five generators create new prompt candidates from the current
  parents; all candidates are scored with the same fixed model and
  validation set, and the best $N$ are kept. Based on Sécheresse et
  al.~\cite{secheresse2025gaapo}.}
  \label{fig:bg-gaapo}
\end{figure}

\begin{bgexample}[Crossover and mutation of two prompts]
\textbf{Parent A:} \emph{``You are a cardiology nurse. Write one short video
scene for patients. Use simple words.''}\\
\textbf{Parent B:} \emph{``Create a storyboard scene in JSON. Explain
abbreviations at first mention.''}\par\smallskip
\textbf{Crossover} (first half of A + second half of B):
\emph{``You are a cardiology nurse. Write one short video scene. Explain
abbreviations at first mention.''}\\
\textbf{Mutation} (strategy ``constraint addition''):
\emph{``\ldots\ Keep the narration under 35 words.''}\par\smallskip
Both children are then scored, and they replace weaker members of the
population only if they achieve a higher score.
\end{bgexample}

% ---------------------------------------------------------------------
\section{Cardiovascular Disease and Patient Education}
\label{sec:bg-medical}
% ---------------------------------------------------------------------

This section introduces the medical application domain of the thesis. It is
written for readers with a computer-science background and only covers the
concepts needed to understand the data and the evaluation.

\subsection{Cardiovascular Disease and Heart Failure}
\label{sec:bg-cvd}

\emph{Cardiovascular diseases} (CVDs) are a group of disorders of the heart
and blood vessels, including coronary heart disease, cerebrovascular
disease and heart failure. According to the World Health Organization, an
estimated 19.8~million people died from CVDs in 2022, which corresponds to
approximately 32\,\% of all global deaths~\cite{who2025cvd}. Many of these
deaths are related to modifiable risk factors such as high blood pressure,
tobacco use, unhealthy diet and physical inactivity, which makes patient
education an important part of prevention and long-term care.

\emph{Heart failure} (HF) is not a single disease but a clinical syndrome
with typical symptoms such as breathlessness, ankle swelling and fatigue. It is
caused by a structural and/or functional abnormality of the heart, which
leads to increased pressures inside the heart and/or an insufficient output of
blood~\cite{mcdonagh2021esc}.
An important measurement is the \emph{left-ventricular ejection fraction}
(LVEF), the percentage of blood that the left ventricle ejects with each
contraction. The European Society of Cardiology (ESC) uses the LVEF to
distinguish three types of heart failure (Table~\ref{tab:bg-hf-types}).

\begin{table}[htbp]
  \centering
  \begin{tabular}{@{}lll@{}}
    \toprule
    \textbf{Abbreviation} & \textbf{Type of heart failure} & \textbf{LVEF} \\
    \midrule
    HFrEF  & HF with reduced ejection fraction       & $\leq 40\,\%$ \\
    HFmrEF & HF with mildly reduced ejection fraction & $41$--$49\,\%$ \\
    HFpEF  & HF with preserved ejection fraction      & $\geq 50\,\%$ \\
    \bottomrule
  \end{tabular}
  \caption{Types of heart failure according to the 2021 ESC
  guidelines~\cite{mcdonagh2021esc}. Additional criteria (symptoms, signs and
  further findings) are required for a diagnosis.}
  \label{tab:bg-hf-types}
\end{table}

Patients with heart failure often have several co-morbidities, such as
hypertension, diabetes or atrial fibrillation, and take several drugs at the
same time. Their questions (``What does an ejection fraction of 25\,\% mean
for me?'') therefore require answers that are correct, personalised and
understandable at the same time, which is the motivating scenario of this
thesis (see the motivating example in Chapter~1).

\subsection{Clinical Practice Guidelines and Recommendation Grading}
\label{sec:bg-guidelines}

\emph{Clinical practice guidelines} are documents in which expert panels
summarise the available evidence and formulate recommendations for diagnosis
and treatment. For heart failure in Europe, the main reference is the 2021
ESC guideline~\cite{mcdonagh2021esc} together with its 2023 focused
update~\cite{mcdonagh2023focused}. Each recommendation is graded along two
dimensions (Table~\ref{tab:bg-grades}): the \emph{class of recommendation}
expresses how strongly an action is recommended, and the \emph{level of
evidence} expresses how strong the underlying scientific evidence is.

\begin{table}[htbp]
  \centering
  \small
  \begin{tabular}{@{}p{2.9cm}p{1.9cm}p{7.6cm}@{}}
    \toprule
    \textbf{Dimension} & \textbf{Grade} & \textbf{Meaning / wording} \\
    \midrule
    Class of        & Class I   & Is recommended or indicated \\
    recommendation  & Class IIa & Should be considered \\
                    & Class IIb & May be considered \\
                    & Class III & Is not recommended \\
    \midrule
    Level of        & Level A & Data from multiple randomised clinical trials or meta-analyses \\
    evidence        & Level B & Data from a single randomised trial or large non-randomised studies \\
                    & Level C & Consensus of experts and/or small studies, retrospective studies, registries \\
    \bottomrule
  \end{tabular}
  \caption{Grading of recommendations in ESC guidelines~\cite{mcdonagh2021esc}.}
  \label{tab:bg-grades}
\end{table}

\begin{bgexample}[A graded recommendation as a Knowledge Graph]
The guideline sentence \emph{``An MRA is recommended for patients with HFrEF
to reduce the risk of HF hospitalisation and death (Class~I, Level~A)''}
(paraphrased from~\cite{mcdonagh2021esc}) can be represented as a small
graph: a node for the recommendation with edges to the drug class
(\textit{MRA}), the patient group (\textit{HFrEF}), the outcome
(\textit{HF hospitalisation and death}), the source document, and a grade
record that has a \textit{classOfRecommendation} (\textit{I}) and a
\textit{levelOfEvidence} (\textit{A}). Keeping the grade and the source as
explicit nodes allows a generated answer to state not only \emph{what} is
recommended, but also \emph{how strongly} and \emph{according to which
guideline}.
\end{bgexample}

\subsection{Health Literacy and Readability}
\label{sec:bg-readability}

Medical correctness alone is not enough; patients must also be able to
understand and use the information.

\begin{bgdefinition}{Health Literacy}
\emph{Health literacy} ``entails people's knowledge, motivation and
competences to access, understand, appraise, and apply health information in
order to make judgments and take decisions in everyday life concerning
healthcare, disease prevention and health promotion to maintain or improve
quality of life during the life course''~\cite{sorensen2012health}.
\end{bgdefinition}

A simple, widely used proxy for how understandable a text is, is a
\emph{readability formula}. The best known is the Flesch Reading Ease
score~\cite{flesch1948new}; Amstad~\cite{amstad1978verstandlich} adapted it to
German, which is the language of the DigiLearnHF platform.

\begin{bgdefinition}{Flesch Reading Ease}
Let $\mathit{ASL}$ be the average sentence length (words per sentence) and
$\mathit{ASW}$ the average number of syllables per word. Then
\begin{align*}
  \mathit{FRE}_{\mathit{en}} &= 206.835 - 1.015 \cdot \mathit{ASL} - 84.6 \cdot \mathit{ASW}
     &&\text{(English,~\cite{flesch1948new})}\\
  \mathit{FRE}_{\mathit{de}} &= 180 - \mathit{ASL} - 58.5 \cdot \mathit{ASW}
     &&\text{(German,~\cite{amstad1978verstandlich})}
\end{align*}
Higher values indicate easier texts; scores of 60--70 correspond to plain
language and scores above 80 to easy texts~\cite{flesch1948new}.
\end{bgdefinition}

\begin{bgexample}[Same content, different readability]
\textbf{Technical:} \emph{``Pharmacological therapy with mineralocorticoid
receptor antagonists reduces hospitalisation in patients with reduced
ejection fraction.''} (14 words, 44 syllables) \\
$\mathit{FRE}_{\mathit{en}} = 206.835 - 1.015\cdot 14 - 84.6\cdot\frac{44}{14} \approx -73$
\;(extremely difficult)\par\smallskip
\textbf{Plain:} \emph{``Some heart medicines lower the risk that you have to
go to the hospital.''} (14 words, 19 syllables) \\
$\mathit{FRE}_{\mathit{en}} = 206.835 - 1.015\cdot 14 - 84.6\cdot\frac{19}{14} \approx 78$
\;(fairly easy)\par\smallskip
Both sentences have the same length; the difference comes entirely from long,
technical words. Readability formulas do not measure correctness, so a
simplified text must still be checked against the original guideline.
\end{bgexample}

\subsection{Educational Videos and Multimedia Learning}
\label{sec:bg-video}

Short educational videos are a popular format for patient education. A
systematic review of video-based interventions for patients with chronic
illnesses found strong evidence that videos improve patients' knowledge of
their condition, whereas effects on clinical outcomes were less
consistent~\cite{deshpande2023video}. The design of such videos can be
guided by the \emph{cognitive theory of multimedia learning}~\cite{mayer2020multimedia},
which assumes that people process information through two channels (visual
and auditory), that each channel has a limited capacity, and that learning
requires active processing. From this theory, Mayer derives design principles
such as:
\begin{itemize}
  \item \textbf{Multimedia principle:} people learn better from words and
        pictures than from words alone.
  \item \textbf{Coherence principle:} extraneous material (decorative
        pictures, unnecessary details) should be excluded.
  \item \textbf{Signalling principle:} key information should be highlighted,
        e.g.\ with short on-screen text.
  \item \textbf{Segmenting principle:} complex content should be presented in
        short, learner-paced segments (scenes).
  \item \textbf{Modality principle:} graphics are better explained with spoken
        narration than with long printed text.
\end{itemize}

Before a video is produced, its content is planned in a \emph{storyboard}.

\begin{bgdefinition}{Storyboard}
A \emph{storyboard} is an ordered sequence of scenes $S = (s_1, \dots, s_n)$
that specifies an educational video before production. Each scene
$s_i = (\mathit{narration}_i, \mathit{visual}_i, \mathit{text}_i)$ consists of
the spoken narration, a description of the visual content, and short
on-screen text that signals the key message.
\end{bgdefinition}

\begin{bgexample}[A one-scene storyboard]
The following storyboard, taken from the dataset used in this thesis, answers
the request \emph{``Create a short patient-education scene that explains the
NYHA classification of heart failure and how it relates to symptom
severity.''}
\begin{lstlisting}[style=bgcode]
{ "title": "Understanding NYHA Heart Failure Classes",
  "scenes": [ {
    "scene_number": 1,
    "narration": "Heart failure is grouped into NYHA Classes I to IV,
                  reflecting how symptoms limit daily activities and
                  guide treatment decisions.",
    "visual_description": "Simple chart showing four icons representing
                           increasing activity limitation.",
    "on_screen_text": "NYHA I-IV: Symptom Levels" } ] }
\end{lstlisting}
The narration carries the explanation (modality principle), the visual
illustrates it (multimedia principle), and the on-screen text highlights the
key term (signalling principle).
\end{bgexample}

A structured storyboard format has two advantages for LLM-based generation:
it can be validated automatically (e.g.\ whether all fields are present and
the narration is short enough), and its medical statements can be checked
scene by scene against the Knowledge Graph before the video is produced.

% ---------------------------------------------------------------------
\section{Summary}
\label{sec:bg-summary}
% ---------------------------------------------------------------------

Figure~\ref{fig:bg-overview} summarises how the concepts introduced in this
chapter relate to each other. Semantic data modelling and RML mappings turn
questionnaires and clinical guidelines into Knowledge Graphs
(Sections~\ref{sec:bg-sparql}--\ref{sec:bg-enrichment}). Rule mining extracts
symbolic regularities from these graphs, which are grounded and verbalised
as chain-of-thought training data for parameter-efficient fine-tuning
(Sections~\ref{sec:bg-peft}, \ref{sec:bg-nesy}). At inference time, the
fine-tuned model receives Knowledge Graph context (Section~\ref{sec:bg-kg-llm})
and a prompt that is optimised automatically (Section~\ref{sec:bg-apo}) to
produce patient-oriented storyboards for educational videos
(Section~\ref{sec:bg-medical}).

\begin{figure}[htbp]
  \centering
  \begin{tikzpicture}[
      font=\footnotesize, >={Stealth[length=2mm]},
      blk/.style={draw, thick, rounded corners=3pt, minimum height=1.05cm,
                  text width=2.15cm, align=center},
      sec/.style={font=\scriptsize\itshape, text=bgGrey, inner sep=2pt}]
    % row 1: symbolic side
    \node[blk, fill=bgGreyLight] (src) at (0,0) {Questionnaires\\\& clinical\\guidelines};
    \node[blk, fill=bgBlueLight, draw=bgBlue] (kg) at (3.85,0) {Knowledge\\Graph (RDF)};
    \node[blk, fill=bgBlueLight, draw=bgBlue] (rules) at (7.7,0) {Symbolic rules\\(AMIE)};
    \node[blk, fill=bgBlueLight, draw=bgBlue] (cot) at (11.55,0) {CoT training\\examples};
    % row 2: neural side
    \node[blk, fill=bgOrangeLight, draw=bgOrange] (ft) at (11.55,-2.9) {LoRA fine-tuned\\LLM};
    \node[blk, fill=bgPurpleLight, draw=bgPurple] (prompt) at (7.7,-2.9) {Optimised\\prompt (GAAPO)};
    \node[blk, fill=bgGreenLight, draw=bgGreen] (story) at (3.85,-2.9) {Storyboard\\script};
    \node[blk, fill=bgGreenLight, draw=bgGreen] (video) at (0,-2.9) {Educational\\video};
    % arrows row 1
    \draw[->, thick] (src) -- node[above, sec] {RML} (kg);
    \draw[->, thick] (kg) -- node[above, sec] {mining} (rules);
    \draw[->, thick] (rules) -- node[above, sec] {grounding} (cot);
    \draw[->, thick] (cot) -- node[right, sec] {fine-tuning} (ft);
    % KG context at inference time
    \draw[->, thick, dashed, bgBlue] (kg.south) -- ++(0,-0.75) -| node[pos=0.25, above, sec, text=bgBlue] {KG context at inference time (RAG)} ([xshift=-6mm]ft.north);
    % arrows row 2
    \draw[->, thick] (prompt) -- node[above, sec] {prompt} (ft);
    \draw[->, thick] (ft.south) -- ++(0,-0.5) -| node[pos=0.25, below, sec] {generation} (story.south);
    \draw[->, thick] (story) -- node[above, sec] {production} (video);
    % layer labels
    \node[anchor=south west, font=\scriptsize\bfseries, text=bgBlue] at ($(src.north west)+(0,0.08)$) {Symbolic knowledge};
    \node[anchor=north west, font=\scriptsize\bfseries, text=bgOrange] at ($(video.south west)+(0,-0.08)$) {Neural generation};
  \end{tikzpicture}
  \caption[Overview of the background concepts]{Overview of how the
  background concepts of this chapter are connected in the
  neuro-symbolic pipeline of this thesis.}
  \label{fig:bg-overview}
\end{figure}
